%----------------------------------------------------------------------------
\chapter{Related work}
\label{chap:RelatedWork}
%----------------------------------------------------------------------------

This chapter introduces the position-based and primal-descent lineages whose
endpoints --- XPBD and VBD --- this thesis compares, and positions my work
relative to the closest prior comparison.

%----------------------------------------------------------------------------
\section{From PBD to XPBD}
\label{sec:PBDtoXPBD}
%----------------------------------------------------------------------------
Position-Based Dynamics (PBD), introduced by M\"uller et~al.~\cite{MULLER2007109},
targets real-time cloth and soft-body simulation in games. Each iteration
projects the predicted positions onto the constraint set --- moving the particles
each constraint couples until the constraint is satisfied --- through
Gauss--Seidel sweeps, and recovers velocities by finite difference. PBD is unconditionally stable, but it has a
well-known flaw: \emph{constraint stiffness depends on the iteration count and the
time step}.

XPBD, introduced by Macklin et~al.~\cite{10.1145/2994258.2994272}, fixes this by
augmenting each constraint with a compliance term $\alpha = 1/k$ and solving for
the associated Lagrange multiplier $\lambda$. The resulting iteration solves an
implicit-Euler problem with arbitrary elastic energies, and constraint stiffness
becomes a true material parameter, independent of the iteration count. As Ma
et~al.~\cite{app142311072} summarize the lineage:

\begin{quote}
``Early research focused on force--object relationships, leading to the
development of models such as the Mass-Spring System [30], Finite Element Method
[31], and Rigid Body Dynamics [32]. However, these models have limitations,
including performance degradation, lack of real-time capability, and unrealistic
motion as simulation complexity increases. PBD was introduced to address these
limitations. Initially developed by Jakobsen [33], PBD provides a stable and
efficient approach by applying wedge-based constraints instead of force-based
methods. Macklin et al. [34] further refined and optimized PBD for real-time
interaction.''
\end{quote}

In short, XPBD was introduced to resolve the fundamental limitation of the
original PBD, in which the stiffness of the material was inherently coupled to the
time-step size and the number of iterations. By introducing a compliant
constraint formulation that incorporates Lagrange multipliers, XPBD allows soft
and stiff materials to be simulated with physical parameters that remain
consistent regardless of the solver's configuration.

\paragraph{Small steps in physics simulation.}
Macklin et~al.~\cite{Macklin2019SmallSI} (the ``Small Steps'' note) later argued
that XPBD converges most reliably with \emph{many short substeps and a single
iteration per substep}, rather than a single long step with many iterations. With only one iteration per
substep the multiplier term entering the update is always zero --- $\lambda$ is
reset each substep and never re-used across iterations --- so there is no need to
store and carry it between iterations. This version is now a standard XPBD configuration, but
for the sake of a fair comparison --- to compare methods across substep and iteration counts --- I deliberately leave this optimization out of
the core comparison.

%----------------------------------------------------------------------------
\section{Vertex Block Descent (VBD)}
\label{sec:VBDLineage}
%----------------------------------------------------------------------------
Vertex Block Descent, due to Chen et~al.~\cite{Chen2024}, treats each vertex as a
tiny optimization block: at each iteration, every vertex moves to the minimizer of
$G$ with all other vertices held fixed. This requires the local $3\times 3$
Hessian of the energies sharing the vertex and a Newton step on that block. Ma
et~al.~\cite{app142311072} place the method as follows:

\begin{quote}
``More recently, the VBD algorithm that has been proposed provides an efficient
approach to parallelize large-scale computations. By processing nodes in block
units, VBD significantly enhances performance, particularly in simulations
involving a high number of nodes [35].''
\end{quote}

The original VBD paper contributes three features beyond the basic block descent.
\emph{Chebyshev acceleration} is an over-relaxation that significantly speeds
convergence, at the cost of the method's energy guarantees. \emph{Adaptive
initialization} seeds each step's initial guess from a blend of the full inertial (\emph{TODO: Write example equations here for adaptive init})
prediction and an inertia-only extrapolation, leaning on the latter for vertices
near rest to avoid over-shooting under gravity. \emph{Graph-coloring parallelism} lets vertices that
share no constraint be updated independently, enabling massive parallelism.

A follow-up, Augmented VBD (AVBD)~\cite{Giles2025}, improves stiffness-ratio
handling, infinite-stiffness handling, and convergence behavior. I restrict the
comparison to the canonical 2024 VBD, since AVBD post-dates my implementation
window; I cite AVBD in \autoref{sec:FlexResults} as the most relevant
counter-evidence on the stiffness-ratio axis.

%----------------------------------------------------------------------------
\section{Prior comparison: Ma et~al.\ (2024)}
\label{sec:Ma2024}
%----------------------------------------------------------------------------
The most directly comparable prior work is Ma et~al.~\cite{app142311072}, who
compare cloth simulation under PBD and VBD in Unity and establish that VBD scales
better than PBD with node count. My work differs in four substantial ways:
\begin{enumerate}
    \item \textbf{XPBD, not PBD.} Plain PBD has the stiffness-coupling flaw that
        XPBD was designed to fix, so a PBD-vs-VBD comparison conflates algorithmic
        differences with this known bug. Using XPBD isolates the primal/dual
        distinction.
    \item \textbf{Canonical VBD.} Ma et~al.'s pseudocode departs in several
        details (initialization, Chebyshev handling) from \cite{Chen2024}. I
        follow canonical VBD, but disable Chebyshev acceleration for the core
        comparison, so that I read off algorithmic properties rather than
        acceleration tricks. Adaptive initialization is kept enabled, since it is
        VBD's inertial prediction and has a direct counterpart in XPBD's explicit
        guess (see \autoref{ssec:DisableChebyshev}).
    \item \textbf{Beyond frame rate.} Ma et~al.\ report fps and node counts. I add
        a Newton ground truth, a solution quality vs.\ computational cost comparison,
        and targeted stress tests (mass ratio, stiffness ratio, large $dt$) keyed to the
        primal/dual distinction.
    \item \textbf{Varying the substep count.} Their comparison kept the substep
        count fixed throughout, so it did not consider how changing it affects each
        algorithm.
\end{enumerate}

In short, this thesis answers a question that Ma et~al.\ explicitly left as future
work: how does VBD compare to \emph{XPBD}, and at \emph{matched quality} rather
than only at matched node counts?

%----------------------------------------------------------------------------
\section{Reference materials}
\label{sec:RefMaterials}
%----------------------------------------------------------------------------
The Physics-Based Simulation book of Li et~al.~\cite{li2026physics} provided
canonical derivations and pedagogical notes; its treatment of implicit Euler as
variational minimization directly informs the framing in \autoref{sec:Variational}.
The \emph{Ten Minute Physics} tutorial series~\cite{tenMinutePhysics} was used as
a sanity-check reference for the correctness of the XPBD implementation.
